% !TeX encoding = UTF-8
% !TeX spellcheck = de_DE
\pdfminorversion=4

%\newcommand{\todo}[1]{\textbf{\color{red}{TODO: #1}}}

\documentclass[german]{beamer} % evtl. FabLab Anpassungen
\usepackage[utf8]{inputenc}
\usepackage[german]{babel}
\usepackage[babel,german=quotes]{csquotes} %\enquote{...} für deutsche Anführungszeichen um ein Wort
\usepackage{color}
\usepackage{amsmath, amssymb, amstext, amsfonts, mathrsfs}
\usepackage{hyperref}
\usepackage{scrhack}
\usepackage{multicol}

\newcommand{\images}{kassenterminal/}
\newcommand{\ccbysa}{\href{https://creativecommons.org/licenses/by-sa/3.0/deed.de}{\includegraphics[height=0.11in]{img/licenses/cc-by-sa.png}}}
\newcommand{\artlibre}{\href{http://artlibre.org/licence/lal/de/}{\includegraphics[height=0.11in]{img/licenses/art_libre.png}}}

\beamertemplatenavigationsymbolsempty

\addtobeamertemplate{footline}{
    \usebeamerfont{footline}
    \usebeamercolor[fg]{footline}
    \hspace{1em}%
    \insertframenumber/\inserttotalframenumber
    \hspace{1em}%
    \insertsectionhead
}{}

%-----------------------------------------------------------------------------------------------------------------

\title{Das Kassenterminal des FAU FabLab}
\author{Patrick Kanzler}
\logo{\includegraphics[width=0.3\linewidth]{img/logo/logo.pdf}}

\begin{document}

%-----------------------------------------------------------------------------------------------------------------

\maketitle

%-----------------------------------------------------------------------------------------------------------------

\frame{\tableofcontents}

%-----------------------------------------------------------------------------------------------------------------

\begin{frame}
\frametitle{Unser Hintergrund}
\begin{itemize}
 \item FabLab als Teil eines Universitäts-Lehrstuhles
 \item Buchhaltung über Lehrstuhl
 \item Materialverkauf und Verkauf von Laserzeit im Lab
 \item früher: offene \enquote{Vertrauenskasse}
   \begin{itemize}
   	\item unübersichtliche Preislisten
   	\item Geldmengen wurden zu groß
    \item hat zu Problemen beim Wechseln geführt
   \end{itemize} 	
\end{itemize}	 
\end{frame}

\begin{frame}
\frametitle{Unsere Anforderungen}
 \begin{itemize}
 	\item Bargeld, später Uni-interne Bezahlkarte
 	\item damals Import der Artikel aus Google-Tabelle, heute ERP
	\item Nachvollziehbarkeit
 \end{itemize}
\end{frame}

\begin{frame}
\frametitle{Unser Kassenterminal}
\begin{columns}[c]
\column{.4\linewidth}
  \only<1>{\includegraphics[width=.9\linewidth]{\images/terminal_geschlossen.jpg}}
  \only<2->{\includegraphics[width=.9\linewidth]{\images/terminal_offen.jpg}}
\column{.59\linewidth}
 \begin{itemize}
 \uncover<1->\item soweit möglich Open-Source-Soft- und Hardware
 \uncover<2->\item Holzkasten, mit Münzzähler, Scheinaufnahme, normaler PC, Touchscreen
 \uncover<3->\item Software: FabLabKasse (\url{https://github.com/fau-fablab/FabLabKasse} )
 \uncover<3->\item Software wird mit git versioniert
 \end{itemize}
 \end{columns}
\end{frame}

\begin{frame}
\frametitle{FabLabKasse}
 \begin{itemize}
  \item geschrieben in Python, GPL
  \item modular gestaltet:
    \begin{itemize}
     \item austauschbares \enquote{Backend} zum Anbinden von Artikeln
     \item Schnittstelle zu Bezahlmöglichkeiten
     \item wir versuchen den Code so allgemein wie möglich zu halten
    \end{itemize}
  \item unser Backend:
    \begin{itemize}
      \item kann Daten aus odoo (ERP) lesen
      \item kann nicht ins ERP schreiben (also auch keine Aktualisierung des Bestandes)
      \item Kassenbuch wird in kassenbuch.py geführt, rudimentäre doppelte Buchführung
    \end{itemize}
  \end{itemize}
\end{frame}
\begin{frame}
\frametitle{FabLabKasse}
  \begin{itemize}
  \item unsere Barzahlung:
     \begin{itemize}
        \item bindet MDB-Münzaufnahme an (Industriestandard)
        \item bindet Geldscheinleser an
        \item günstige Geräte haben oftmals Bugs in Firmware
        \item[$\rightarrow$] \enquote{paranoider} Code, um Fehler zu erkennen
     \end{itemize}
  \item Zusatzfeature:
     \begin{itemize}
       \item User können mit unserer App Warenkörbe zusammen stellen
       \item dieser Warenkorb kann an das Kassenterminal übertragen werden
     \end{itemize} 
 \end{itemize}
\end{frame}

\begin{frame}
\frametitle{Was ist noch bei uns zu tun?}
 \begin{itemize}
   \item richtige Anbindung ERP (Bestand, Verkäufe ins ERP, ...)
   \item Verallgemeinerung des Kassenbuch und zugehörigem Backend
   \item Anbindung der Uni-internen Bezahlkarten
   \item viele kleine Umbauten, um den Code allgemeingültiger zu machen
 \end{itemize} 
\end{frame}

\begin{frame}
\frametitle{Was ist zu tun, um den Code bei dir einzusetzen?}
  \begin{itemize}
    \item Dies ist keine Plug-and-Play-Lösung! Vielmehr auf uns zugeschnitten
    \item Möglicher Pfad:
      \begin{itemize}
         \item Hardware zusammensuchen (wobei auch eine offene Kasse ohne \enquote{Automatisierung} möglich wäre)
         \item Treiber ergänzen (falls nötig)
         \item Backend schreiben
         \item bei bestehendem ERP o.\,ä. daran anpassen
         \item oder z.\,B. unser Kassenbuch anpassen
      \end{itemize}
    \item wir würden mit Rat zur Seite stehen, haben aber nicht die Kapazität, diese Anpassungen selbst durchzuführen
  \end{itemize}
\end{frame}

\begin{frame}
\frametitle{Kontakt}
 \begin{itemize}
  \item \url{kontakt@fablab.fau.de}
  \item oder z.\,B. Issue auf github ( \url{https://github.com/fau-fablab/FabLabKasse} )
  \item oder \url{patrick.kanzler@fablab.fau.de} (ich stelle gerne den weiteren Kontakt her)
 \end{itemize}
\end{frame}


\end{document}
